\documentclass[12pt]{article}
\usepackage{amsmath,amssymb}
\usepackage{booktabs}
\usepackage{multirow}
\usepackage{natbib}
\usepackage{hyperref}
\usepackage{microtype}
\usepackage{xcolor}
\usepackage{mdframed}
\usepackage{enumitem}
\usepackage{graphicx}
\usepackage{caption}
\usepackage{geometry}
\geometry{margin=1in}
\usepackage{xr}
\externaldocument{hrf}

% Referee comment box
\newmdenv[
  linecolor=gray!50,
  backgroundcolor=gray!8,
  innerleftmargin=10pt,
  innerrightmargin=10pt,
  innertopmargin=8pt,
  innerbottommargin=8pt,
  skipabove=10pt,
  skipbelow=4pt
]{refbox}

\newcommand{\comment}[1]{\begin{refbox}\itshape #1\end{refbox}}
\newcommand{\response}{\medskip\noindent\textbf{Response.}\ }
\newcommand{\hrf}{\textsc{hrf}}
\newcommand{\rf}{\textsc{rf}}
\newcommand{\todo}[1]{\par\medskip\noindent\fcolorbox{red}{red!8}{\parbox{\dimexpr\linewidth-2\fboxsep-2\fboxrule}{\textbf{\color{red}TODO (internal --- remove before submission):}\ \color{red}#1}}\par}

\title{Response to the Associate Editor\\[6pt]
  \large ``The Hedged Random Forest''\\
  \normalsize Journal of Financial Econometrics --- JFEC-2025-167}
\author{E.\ Beck \and D.\ Kozbur \and M.\ Wolf}
\date{}

\begin{document}
\maketitle
\thispagestyle{empty}

We thank the Associate Editor for the encouraging assessment of the paper's
potential and for the constructive guidance on how to strengthen it. Below we
address each of the three points in turn. All section, table, and figure
references are to the revised manuscript.

% ============================================================
\section*{Comment 1: A financial-econometric application}
% ============================================================

\comment{The empirical section would benefit from a closer match with the
readership and standards of Journal of Financial Econometrics. The current
evidence is useful in showing that the proposed hedged random forest improves
upon the standard random forest and several weighted-random-forest variants
on benchmark regression datasets. However, it does not yet fully demonstrate
how the method performs in a financial-econometric setting. I would
encourage the authors to add one substantive financial-econometric
application, such as return prediction, volatility or realized-volatility
forecasting, option-implied volatility forecasting, credit-risk prediction,
macro-financial forecasting, or risk-premium estimation. The application
should be framed as a genuine out-of-sample forecasting exercise, with
rolling or expanding-window evaluation where appropriate. At the same time,
I do not think the paper needs to become a broad horse-race comparison
against all modern tree-based algorithms.}

\response
We have added a new application, Section~\ref{sec:fredmd}, that directly
addresses this comment. Rather than a single hand-picked series, we forecast
a broad panel of 41 target series drawn from two groups of the FRED-MD
database \citep{mccracken:ng:2016}: interest-rate levels, credit and term
spreads, and exchange rates (the \emph{Interest and Exchange Rates} group),
and price indices (the \emph{Prices} group). Scanning across an entire class
of targets, rather than reporting a single series, was a deliberate choice:
it lets us show \emph{where} the hedging mechanism helps within a coherent
financial-econometric setting, without turning the exercise into a
horse-race across competing algorithms, which you correctly note the paper
does not need.

The exercise is a genuine expanding out-of-sample forecasting exercise in
the sense you describe: for each target and each of 432 monthly origins
(1990-05 through 2026-04), we fit the random forest (and \hrf) on a rolling
360-month window, forecast the next month, and record the realized forecast
error. We adapt the time-series estimators of $\hat\mu$ and $\hat\Sigma$
developed in our companion paper on inflation forecasting
\citep{beck:wolf:hrf-inf} (EWMA estimation combined with linear shrinkage)
rather than the off-the-shelf estimators used for the i.i.d.\ benchmark
datasets in Section~\ref{sec:results}, since --- as we discuss in
Section~\ref{sec:fredmd} --- suitable estimators of the two hedging inputs
depend on the nature and frequency of the data for time-series applications.

Table~\ref{table:fredmd-summary} and Figure~\ref{fig:fredmd-boxplot}
summarize the results by group, with the full 41-series breakdown in
Appendix~B, Table~13. Across all 41 series, \hrf\ achieves a lower RMSE than
\rf\ for 32 series (78\%), with a mean RMSE ratio of 0.982. The gains are
strongest and most uniform for the eight credit and term spreads (\hrf\
wins for every single series, mean ratio 0.938), more modest for the price
indices (mean ratio 0.991), mixed for interest-rate levels (mean ratio
0.995, concentrated at the short end of the curve), and essentially absent
for exchange rates (mean ratio 1.000), consistent with the well-documented
difficulty of beating a random walk in exchange-rate forecasting. This
last point is itself informative: it shows that \hrf\ does not manufacture
spurious gains where none should be expected.

% ============================================================
\section*{Comment 2: Incremental value relative to modern tree-based methods}
% ============================================================

\comment{Refining the contribution is important because unequal and signed
contributions of trees are already familiar in boosted-tree methods. The
authors should therefore show whether HRF offers incremental value relative
to modern tree-based machine-learning tools, or clarify that its
contribution is instead a low-tuning post-estimation improvement to random
forests. A carefully designed financial-econometric application, together
with some diagnostic evidence on when the weighting scheme works, would
make the empirical contribution considerably more convincing.}

\response
We take the second of the two options you outline: we clarify, rather than
attempt to out-perform, boosted-tree methods, for the following reason.
Gradient boosting \citep{friedman:2001} and its variants construct an
entirely different ensemble than the random forest: trees are fit
sequentially to the residuals of the current ensemble, and performance
typically depends on the joint tuning of several hyperparameters (the
learning rate, tree depth, number of boosting rounds, and regularization).
\hrf, by contrast, is a \emph{post-estimation} re-weighting of an
already-trained, standard random forest. It requires no change to how the
individual trees are grown and no retraining, and it has a single
user-facing hyperparameter, the leverage constraint $\kappa$, for which we
recommend --- and use throughout the paper, including in the new
Section~\ref{sec:fredmd} --- the default choice $\kappa = 2$ (see
Section~\ref{sec:hrf}). We therefore view \hrf\ as a low-tuning,
``point-and-shoot'' improvement (this is literally how we describe the
i.i.d.\ estimators in Section~\ref{sec:fredmd}) that a random-forest user
can apply on top of an existing forest, rather than as a competitor to
boosting-type ensembles built from the ground up. We agree that a direct
horse-race against boosting would be a natural and interesting extension,
but, per your own comment above, we have deliberately kept the paper's
empirical scope to \hrf\ versus \rf\ and previously proposed weighted
random forests, rather than broadening it into a comparison against all
modern tree-based algorithms. This is a deliberate scoping strategy on our
part, not an oversight: \hrf\ is positioned throughout the paper as a
post-estimation refinement of a standard random forest, and we view a fair,
properly tuned comparison against boosted-tree methods as a separate
undertaking that deserves its own paper rather than a section here.

\todo{We are about to tell the AE that we will \emph{not} run a
head-to-head comparison against boosting or other modern tree ensembles in
this paper --- framing \hrf\ purely as a post-estimation add-on and pushing
a boosting comparison to future work. This is a real strategic commitment,
and it isn't settled on our end yet: we should all agree before this goes
out. If we'd rather leave the door open (e.g.\ a referee could ask for this
in a later round, or we may want it for a follow-up paper), let's discuss
before submitting.}

On the diagnostic evidence you request: the FRED-MD results just described
already provide one clear pattern. The gains from \hrf\ are largest and
most uniform for the credit and term spreads, a group in which the
underlying predictors (default risk, liquidity conditions, and the level of
rates) are heterogeneous and only loosely related, plausibly giving rise to
correspondingly heterogeneous, less mutually redundant tree forecasts for
\hrf\ to exploit; the gains are essentially absent for exchange rates, which
are close to unpredictable by any method. A second, independent piece of
diagnostic evidence was reported in our response to Reviewer 1 (Additional
Comment 1): across the 17 i.i.d.\ benchmark datasets of
Section~\ref{sec:results}, the dispersion in individual trees' out-of-bag
accuracy is the one dataset characteristic (of five considered) that is
significantly related to \hrf's performance gain (Spearman $\rho = -0.72$,
$p = 0.002$), with datasets with the most heterogeneous trees showing the
largest gains. Both pieces of evidence point in the same direction: \hrf\
helps most when trees are not all equally reliable, which is exactly the
situation its weighting scheme is designed to exploit.

% ============================================================
\section*{Comment 3: Similarity score and prior circulation}
% ============================================================

\comment{Lastly, the similarity check reports a similarity score of 42\%,
which has raised some concern at the editorial office. Having reviewed the
report, I understand that much of the overlap appears to be with an earlier
working-paper version of the same project. The authors would benefit from
further revising the writing, presentation, and framing of the manuscript
so that the submitted version is more clearly differentiated from the
earlier version, while ensuring that any prior circulation is properly
acknowledged.}

\response
You are correct: an earlier version of this paper circulated under the same
title as a working paper, \citet{beck:kozbur:wolf:2024}. This JFEC
submission is the first formal journal submission of that project; the
working paper itself was never submitted elsewhere. We had not yet added an
explicit acknowledgment of the prior circulation to the manuscript and are
glad to add a title-page footnote citing the working paper if the editorial
office considers this necessary --- we did not want to make further changes
to the main text before confirming the exact wording and placement you and
the editorial office would prefer.

On differentiation from the working paper: relative to that version, the
manuscript has been substantially revised, both in response to Reviewer~1
and in response to this letter. Reviewer~1's comments led us to add two
additional benchmark comparisons not present in the working paper --- a
ridge combination scheme and a bias-free minimum-variance variant,
HRF-MinVar (Section~\ref{sec:comp1}) --- a discussion of negative weights
in relation to tree correlations (Section~\ref{sec:neg-weights}), a remark
clarifying that the leverage constraint yields dense rather than sparse
solutions (Remark~\ref*{rem:dense}), and a unification of notation
throughout the paper around the shrinkage/sample $\times$ $\kappa \in
\{1,2\}$ labeling scheme. This letter adds a further, substantial change:
the new financial-econometric application of Section~\ref{sec:fredmd},
which has no counterpart in the working paper. Between the two rounds of
revision, a large share of the paper's empirical content and exposition is
new relative to the circulated working paper, and we would welcome any
further guidance from the editorial office on presentation changes that
would help differentiate the two versions.

% ============================================================
\bibliographystyle{plainnat}
\bibliography{wolf}
% ============================================================

\end{document}
